\chapter{\eh{rocket} Protocolos Modernos y Tecnologías Google}

\begin{concepto}
La web de hoy es radicalmente distinta a la de los años 90. Los protocolos
que estudiamos en capítulos anteriores (TCP, HTTP/1.1) tienen décadas de
antigüedad y muestran su edad bajo la carga del internet moderno.

Este capítulo cubre la próxima generación: protocolos diseñados para
redes móviles, streaming de video en 4K, videollamadas en tiempo real,
y los sistemas distribuidos masivos de Google, Meta y Cloudflare.
\end{concepto}

\section{\eh{zap} QUIC --- Quick UDP Internet Connections (RFC 9000)}

\begin{concepto}
\textbf{QUIC} es un protocolo de transporte sobre \textbf{UDP} que combina
lo mejor de TCP + TLS en un solo handshake. Desarrollado por Google en 2012,
estandarizado por IETF en 2021 (RFC 9000).

\emoji{light-bulb} \textbf{Analogía:} TCP es como llamar a un banco
--- tienes que autenticarte, esperar verificación, y solo entonces
hablar. QUIC es como un banco con ``acceso rápido'' donde te reconocen
de visitas anteriores y entras directo al grano.
\end{concepto}

\subsection{El problema que QUIC resuelve}

TCP tiene problemas fundamentales con el internet moderno:

\begin{itemize}
    \item \textbf{Handshake lento:} TCP (1 RTT) + TLS 1.2 (2 RTT) = 3 RTTs antes
          del primer byte de datos. En móviles con 100ms de latencia: 300ms
          de espera solo para empezar.
    \item \textbf{HOL Blocking a nivel TCP:} HTTP/2 resolvió el HOL a nivel
          de aplicación con multiplexing, pero si un paquete TCP se pierde,
          \textit{todos} los streams quedan bloqueados en el buffer.
    \item \textbf{Ossification:} middleboxes (NATs, firewalls) han ``osificado''
          TCP --- no puedes cambiar el protocolo TCP porque millones de
          dispositivos en el mundo asumen comportamiento específico.
\end{itemize}

\subsection{Solución QUIC}

\begin{lstlisting}
  TCP + TLS 1.2:
    ├── TCP SYN / SYN-ACK / ACK       (1 RTT)
    ├── TLS ClientHello / ServerHello  (1 RTT)
    ├── TLS Certificate / Finished     (1 RTT)
    └── Primera petición HTTP          (1 RTT) = 4 RTTs en total

  QUIC (primera conexión):
    ├── QUIC Initial (transport + TLS crypto combinados)  (1 RTT)
    └── Primera petición HTTP/3                           (1 RTT) = 2 RTTs

  QUIC (conexión repetida con 0-RTT):
    └── Request + datos inmediatamente                   = 0 RTTs extra
\end{lstlisting}

\begin{table}[h!]
\centering
\renewcommand{\arraystretch}{1.4}
\begin{tabular}{lll}
\toprule
\textbf{Característica} & \textbf{TCP + TLS} & \textbf{QUIC} \\
\midrule
Protocolo base      & IP                    & UDP \\
Handshake           & 3 RTTs (TCP+TLS)      & 1 RTT (combinado) \\
0-RTT resumption    & \emoji{cross-mark} No & \emoji{check-mark-button} Sí \\
HOL Blocking        & Nivel TCP y app       & Solo nivel stream (no cross-stream) \\
Cifrado             & TLS separado          & TLS 1.3 integrado, obligatorio \\
Migración de red    & Reinicia conexión     & \emoji{check-mark-button} Connection ID, sobrevive \\
\bottomrule
\end{tabular}
\caption{TCP+TLS vs QUIC}
\end{table}

\begin{moderno}
\textbf{\emoji{rocket} Connection Migration:} si cambias de Wi-Fi a 4G mientras
descargas un archivo, TCP reinicia la conexión (el socket es IP:puerto). QUIC usa
un \textbf{Connection ID} opaco, independiente de la IP. Tu descarga continúa
sin interrupciones. YouTube en Android lleva usando esto desde 2013.
\end{moderno}

\begin{cuidado}
\textbf{Por qué UDP y no un nuevo protocolo?} Los routers y firewalls están
programados en hardware para TCP/UDP --- añadir un Layer 4 nuevo tomaría
décadas en desplegarse. UDP actúa como contenedor neutral que el hardware
ya conoce; QUIC construye todas las garantías encima.
\end{cuidado}

\section{\eh{globe-with-meridians} HTTP/3 (RFC 9114)}

\begin{concepto}
\textbf{HTTP/3} = HTTP/2 semántica + QUIC como transporte. No es solo
``HTTP sobre UDP'' --- es una reescritura completa del stack.

Adoptado por: Google (100\% del tráfico de YouTube), Cloudflare, Meta,
Shopify. Ya representa \textbf{$\sim$30\% del tráfico web global} (2024).
\end{concepto}

\begin{table}[h!]
\centering
\renewcommand{\arraystretch}{1.3}
\begin{tabular}{llll}
\toprule
\textbf{Versión} & \textbf{Año} & \textbf{Transporte} & \textbf{Clave} \\
\midrule
HTTP/1.1 & 1999 & TCP  & Texto plano, pipeline con HOL \\
HTTP/2   & 2015 & TCP  & Multiplexing binario, HPACK, server push \\
HTTP/3   & 2022 & QUIC & Sin HOL TCP, 0-RTT, connection migration \\
\bottomrule
\end{tabular}
\caption{Evolución de HTTP}
\end{table}

\subsection{QPACK --- Compresión de headers en HTTP/3}

HTTP/2 usa HPACK para comprimir headers --- funciona bien sobre TCP donde
el orden está garantizado. QUIC entrega paquetes fuera de orden (streams
independientes), así que HPACK rompería si un paquete de tabla llega tarde.

\textbf{QPACK} (RFC 9204) resuelve esto: permite compresión de headers
sin bloqueo inter-stream. Usa una tabla estática predefinida y una tabla
dinámica con control de riesgo de bloqueo.

\section{\eh{satellite-antenna} gRPC --- Google Remote Procedure Call}

\begin{concepto}
\textbf{gRPC} es un framework de \textit{Remote Procedure Call} (RPC) de alto
rendimiento. Desarrollado por Google (2015, open source). Sobre HTTP/2 con
Protocol Buffers para serialización.

\emoji{light-bulb} \textbf{Analogía:} REST es como mandar un email con
instrucciones en JSON (texto, lento de parsear). gRPC es como hablar
directamente por teléfono en un idioma optimizado: binario, rápido,
con definición de tipos estricta.
\end{concepto}

\subsection{REST vs gRPC}

\begin{table}[h!]
\centering
\renewcommand{\arraystretch}{1.4}
\begin{tabular}{lll}
\toprule
\textbf{Característica} & \textbf{REST/JSON} & \textbf{gRPC/Protobuf} \\
\midrule
Serialización   & JSON texto plano      & Protobuf binario \\
Velocidad parse & Lento                 & 5--10x más rápido \\
Tamaño payload  & Mediano               & 3--10x más pequeño \\
Tipado          & Ninguno en runtime    & Estricto (schema .proto) \\
Streaming       & Manual (SSE/WebSocket)& Nativo (server/client/bidi) \\
HTTP version    & HTTP/1.1 o HTTP/2     & HTTP/2 (requerido) \\
Browser native  & \emoji{check-mark-button} Sí & \emoji{cross-mark} No (necesita grpc-web) \\
Uso             & APIs públicas         & Microservicios internos \\
\bottomrule
\end{tabular}
\caption{REST vs gRPC}
\end{table}

\subsection{Tipos de streaming en gRPC}

\begin{lstlisting}
  1. Unary (request-response normal):
     cliente ──── Request ────► servidor
     cliente ◄─── Response ─── servidor

  2. Server Streaming (servidor manda muchos mensajes):
     cliente ──── Request ────► servidor
     cliente ◄─── Message 1 ── servidor
     cliente ◄─── Message 2 ── servidor
     (útil: stock ticker, logs en vivo)

  3. Client Streaming (cliente manda muchos mensajes):
     cliente ──── Chunk 1 ────► servidor
     cliente ──── Chunk 2 ────► servidor
     cliente ◄─── Response ─── servidor
     (útil: upload de archivos grandes)

  4. Bidirectional Streaming:
     cliente ◄──► servidor
     (útil: chat, videojuegos en tiempo real)
\end{lstlisting}

\begin{moderno}
\textbf{\emoji{rocket} gRPC en Google:} internamente, Google tiene más de
10 billones de RPCs por segundo usando gRPC/Stubby entre sus microservicios.
Google Maps, Google Photos, Gmail: todos se comunican internamente con este
protocolo. YouTube usa gRPC para coordinar entre sus 100+ microservicios.
\end{moderno}

\begin{codigo}
\begin{lstlisting}
// Archivo .proto: define la API como código
syntax = "proto3";

service UsuarioService {
    rpc GetUsuario(UsuarioRequest) returns (UsuarioResponse);
    rpc StreamNoticias(Empty) returns (stream Noticia);
}

message UsuarioRequest { int32 id = 1; }
message UsuarioResponse {
    int32 id = 1;
    string nombre = 2;
    string email = 3;
}
\end{lstlisting}
\end{codigo}

\section{\eh{racing-car} BBR --- Bottleneck Bandwidth and RTT}

\begin{concepto}
\textbf{BBR} es un algoritmo de control de congestión TCP desarrollado por
Google (2016). Reemplaza CUBIC en los servidores de Google. La diferencia
es filosófica: CUBIC reacciona a pérdida de paquetes; BBR modela físicamente
la red.
\end{concepto}

\subsection{El problema con CUBIC}

CUBIC y su predecesor RENO asumen que pérdida de paquete = congestión.
Esto funcionaba en redes cableadas de los 90. Problemas modernos:

\begin{itemize}
    \item \textbf{Buferbloat:} routers modernos tienen buffers enormes.
          Cuando el buffer está lleno, el RTT se dispara pero no hay pérdida
          --- CUBIC sigue enviando rápido, y la latencia es terrible.
    \item \textbf{WiFi y LTE:} pierden paquetes por interferencia (no congestión)
          --- CUBIC frena innecesariamente.
    \item \textbf{Links transcontinentales:} alta capacidad pero alta latencia,
          subexplotados por CUBIC.
\end{itemize}

\subsection{Cómo funciona BBR}

BBR modela dos parámetros físicos de la red:
\begin{enumerate}
    \item \textbf{BtlBw (Bottleneck Bandwidth):} máxima tasa de entrega observada
          --- qué tan rápido puede la red en realidad?
    \item \textbf{RTprop (Round-Trip propagation time):} RTT mínimo observado
          --- cuánto tarda la señal, sin congestión?
\end{enumerate}

\begin{lstlisting}
  Operación ideal = BtlBw × RTprop
  (Ecuación de bandwidth-delay product)

  Si envías MÁS rápido que BtlBw  → llenas buffers → mayor RTT → latencia
  Si envías MÁS lento que BtlBw   → desperdicias capacidad

  BBR se mantiene justo en el punto óptimo, midiendo constantemente.
\end{lstlisting}

\begin{moderno}
\textbf{\emoji{rocket} Resultados reales de BBR en Google:}
\begin{itemize}
    \item Google.com: \textbf{4\% más rápido} globalmente
    \item YouTube: \textbf{14\% menos buffering} en el mundo
    \item Brasil (alta latencia): \textbf{6x más throughput} con BBR vs CUBIC
    \item Con 1\% de pérdida aleatoria: BBR tiene \textbf{800x más
          throughput} que CUBIC
\end{itemize}
BBR v3 (2023) está llegando a Linux kernel 6.x y Android.
QUIC usa BBR por defecto.
\end{moderno}

\section{\eh{telephone-receiver} WebRTC --- Web Real-Time Communication}

\begin{concepto}
\textbf{WebRTC} (RFC 8825) permite comunicación peer-to-peer de audio,
video y datos directamente en el navegador, sin plugins. Google lo compró
(GlobalIPSolutions, 2010) y lo donó al W3C.

Es lo que usan: Google Meet, Discord, Facebook Messenger, Jitsi.
\end{concepto}

\subsection{Stack de protocolos WebRTC}

\begin{lstlisting}
  ┌──────────────────────────────────────────────┐
  │              Tu aplicación                    │
  ├──────────────────────────────────────────────┤
  │  Audio/Video: Opus + VP8/VP9/AV1/H.264        │
  ├──────────────────────────────────────────────┤
  │    SRTP --- Secure Real-time Transport        │
  ├──────────────────────────────────────────────┤
  │    DTLS --- Datagram TLS (TLS sobre UDP)      │
  ├──────────────────────────────────────────────┤
  │    ICE --- Interactive Connectivity           │
  │    (coordina STUN y TURN para NAT traversal)  │
  ├──────────────────────────────────────────────┤
  │            UDP (o TCP como fallback)          │
  └──────────────────────────────────────────────┘
\end{lstlisting}

\subsection{ICE, STUN y TURN --- NAT Traversal}

La mayoría de dispositivos están detrás de NAT. Para una videollamada
peer-to-peer, los dos peers deben encontrarse, pero ninguno conoce la
IP pública del otro.

\begin{table}[h!]
\centering
\renewcommand{\arraystretch}{1.4}
\begin{tabular}{lll}
\toprule
\textbf{Protocolo} & \textbf{Función} & \textbf{Cuándo} \\
\midrule
STUN & Descubre tu IP pública y puerto NAT & Siempre, primero \\
TURN & Relay de tráfico por servidor intermedio & Cuando P2P falla \\
ICE  & Prueba todos los candidatos y elige el mejor & Coordina todo \\
\bottomrule
\end{tabular}
\caption{ICE/STUN/TURN en WebRTC}
\end{table}

\begin{ejemplo}
\textbf{Por qué DTLS y no TLS?} TLS requiere TCP (streams ordenados).
WebRTC usa UDP para baja latencia en video. DTLS es TLS adaptado para UDP:
maneja retransmisiones de handshake, pero una vez establecido, el
cifrado funciona packet-by-packet sin esperar orden.
Audio/video perdido es aceptable; una llamada con 200ms de latencia, no.
\end{ejemplo}

\section{\eh{magnifying-glass-tilted-left} DoH y DoQ --- DNS moderno}

\begin{concepto}
DNS clásico viaja en texto claro sobre UDP puerto 53. Cualquiera en tu
red (tu ISP, un café con Wi-Fi malicioso) puede ver todos tus
queries DNS: cada sitio que intentas visitar, antes de que cargue.
\end{concepto}

\begin{table}[h!]
\centering
\renewcommand{\arraystretch}{1.4}
\begin{tabular}{llll}
\toprule
\textbf{Protocolo} & \textbf{Puerto} & \textbf{Sobre} & \textbf{Adoptado por} \\
\midrule
DNS clásico         & 53 (UDP/TCP) & Texto claro & Todos (legacy) \\
DoT (DNS over TLS)  & 853          & TCP + TLS   & Android 9+, iOS \\
DoH (DNS over HTTPS)& 443          & HTTPS       & Firefox, Chrome, Cloudflare \\
DoQ (DNS over QUIC) & 853          & QUIC        & En adopción (RFC 9250) \\
\bottomrule
\end{tabular}
\caption{Evolución del DNS cifrado}
\end{table}

\begin{moderno}
\textbf{\emoji{rocket} Cloudflare 1.1.1.1:} resolver DNS público más
rápido del mundo (lanzado 2018). Usa DoH y DoT. Promete no registrar
IPs de usuarios. Tiene \textbf{1.1.1.2} con filtrado de malware y
\textbf{1.1.1.3} con filtrado adicional. Google ofrece \textbf{8.8.8.8}
con las mismas capacidades. Ambos tienen latencia $<$10ms desde México.
\end{moderno}

\begin{cuidado}
\textbf{DoH y privacidad:} DoH oculta tus queries de la red local, pero
el resolver (Cloudflare, Google) sí ve todo. Cambias de que te espíe
tu ISP a que te espíe Cloudflare. Evalúa en qué entidad confías más.
\end{cuidado}

\section{\eh{eye} ECH --- Encrypted Client Hello}

\begin{concepto}
Con HTTPS el contenido está cifrado. Pero el \textbf{Server Name
Indication (SNI)} --- el nombre del dominio que visitas --- viaja en
\textbf{texto claro} en el ClientHello de TLS. Tu ISP o un firewall
pueden saber exactamente qué sitio visitas aunque uses HTTPS.

\textbf{ECH} (Encrypted Client Hello) cifra el SNI usando la clave
pública del servidor, publicada en DNS (registro HTTPS).
\end{concepto}

\begin{lstlisting}
  Sin ECH:
  ClientHello {
    server_name: "ejemplo-bloqueado.com"   <- VISIBLE en texto claro
    ...datos cifrados...
  }

  Con ECH:
  ClientHello (outer) {
    server_name: "cloudflare-esni.com"     <- dominio genérico (visible)
    encrypted_client_hello: [              <- inner ClientHello cifrado
      server_name: "ejemplo.com"           <- real, CIFRADO
      ...
    ]
  }
\end{lstlisting}

\begin{moderno}
\textbf{\emoji{rocket} Adopción:} Firefox habilitó ECH por defecto en 2023.
Chrome lo tiene desde 2024. Cloudflare soporta ECH para todos sus dominios.
Cuando ECH esté universalmente adoptado, los firewalls censores por SNI
serán inefectivos sin interceptar la clave privada del servidor.
\end{moderno}

\section{\eh{satellite-antenna} Wi-Fi 7 --- IEEE 802.11be}

\begin{concepto}
\textbf{Wi-Fi 7} (802.11be) fue aprobado como estándar en 2024. La mejora
más revolucionaria: \textbf{Multi-Link Operation (MLO)}.
\end{concepto}

\begin{table}[h!]
\centering
\renewcommand{\arraystretch}{1.3}
\begin{tabular}{lllll}
\toprule
\textbf{Estándar} & \textbf{Wi-Fi} & \textbf{Año} & \textbf{Max teórico} & \textbf{Bandas} \\
\midrule
802.11n  & Wi-Fi 4   & 2009      & 600 Mbps          & 2.4 + 5 GHz \\
802.11ac & Wi-Fi 5   & 2013      & 3.5 Gbps          & 5 GHz \\
802.11ax & Wi-Fi 6/6E & 2019/2021 & 9.6 Gbps         & 2.4+5+6 GHz \\
802.11be & Wi-Fi 7   & 2024      & \textbf{46 Gbps}  & 2.4+5+6 GHz ampliado \\
\bottomrule
\end{tabular}
\caption{Evolución Wi-Fi}
\end{table}

\subsection{Multi-Link Operation (MLO)}

\begin{itemize}
    \item \textbf{Wi-Fi 4/5/6:} tu dispositivo conecta a \textit{una} banda a la vez
          (2.4 GHz O 5 GHz). Si hay interferencia, migras entre bandas.
    \item \textbf{Wi-Fi 7 con MLO:} conecta a \textit{múltiples bandas simultáneamente}
          (2.4 + 5 + 6 GHz en paralelo real).
\end{itemize}

\begin{lstlisting}
  MLO en acción:

  Dispositivo Wi-Fi 7            Router Wi-Fi 7
    link 2.4 GHz ─────────────────────────►
    link 5.0 GHz ─────────────────────────►
    link 6.0 GHz ─────────────────────────►

  El sistema elige dinámicamente qué paquete va por qué link
  según congestión e interferencia. Si 5 GHz tiene interferencia,
  automáticamente usa 6 GHz sin interrupción de la conexión.
\end{lstlisting}

Beneficios de MLO:
\begin{itemize}
    \item \textbf{Latencia reducida:} datos urgentes van por el link más libre
    \item \textbf{Confiabilidad:} si un link falla, los otros continúan
    \item \textbf{Mayor throughput:} agregación real de múltiples links
\end{itemize}

\begin{moderno}
\textbf{\emoji{rocket} 320 MHz de ancho de banda:} Wi-Fi 6 usa hasta 160 MHz.
Wi-Fi 7 usa hasta \textbf{320 MHz} (solo en 6 GHz). Con 4096-QAM y
16 antenas MIMO, alcanza 46 Gbps teóricos. En la práctica, espera
5--10 Gbps en condiciones reales --- suficiente para decenas de streams
de 8K simultáneos.
\end{moderno}

\section{\eh{shield} WireGuard --- VPN Moderna}

\begin{concepto}
\textbf{WireGuard} (2015) es un protocolo VPN extremadamente simple y
rápido. Comparado con OpenVPN (2001) y IPsec: código mucho más pequeño,
criptografía moderna obligatoria, y velocidad cercana a wireline.
\end{concepto}

\begin{table}[h!]
\centering
\renewcommand{\arraystretch}{1.3}
\begin{tabular}{llll}
\toprule
\textbf{Característica} & \textbf{OpenVPN} & \textbf{IPsec} & \textbf{WireGuard} \\
\midrule
Líneas de código    & $\sim$70,000   & $\sim$400,000  & \textbf{$\sim$4,000} \\
Protocolo base      & TCP o UDP      & IP directo     & UDP \\
Cripto obligatoria  & Configurable   & Configurable   & \textbf{Fija, moderna} \\
Handshake           & Lento (TLS)    & IKEv2          & Rápido (1-RTT) \\
Kernel integration  & Userspace      & Parcial        & \textbf{Linux kernel 5.6+} \\
\bottomrule
\end{tabular}
\caption{Comparativa VPN}
\end{table}

WireGuard usa exclusivamente:
\begin{itemize}
    \item \textbf{ChaCha20-Poly1305:} cifrado simétrico (rápido en móviles sin AES-NI)
    \item \textbf{Curve25519:} Diffie-Hellman para intercambio de claves
    \item \textbf{BLAKE2s:} función hash
    \item \textbf{SipHash24:} para la tabla de routing interna
\end{itemize}

No hay negociación de cifrado --- elimina ataques de downgrade de criptografía.

\begin{moderno}
\textbf{\emoji{rocket} WireGuard en producción:} integrado en Linux kernel 5.6
(2020), Android 12+, iOS 15+. Cloudflare usa WireGuard en su servicio
WARP (1.1.1.1 VPN). En benchmarks, WireGuard alcanza velocidades
\textbf{3--4x más rápidas} que OpenVPN en el mismo hardware.
\end{moderno}

\section{\eh{office-building} Zero Trust y BeyondCorp}

\begin{concepto}
\textbf{Zero Trust} es un modelo de seguridad que asume que \textit{ninguna
red es confiable} --- ni la interna. Contrasta con el modelo tradicional
de ``fortaleza y foso'': todo dentro del firewall es confiable.

\emoji{light-bulb} \textbf{Analogía:} modelo tradicional = castillo con
muralla. Si cruzas el firewall, eres confiable y puedes ir a cualquier
habitación. Zero Trust = edificio moderno con tarjetas de acceso en
\textit{cada puerta}. Aunque entres al edificio, cada habitación verifica
tu identidad.
\end{concepto}

\subsection{BeyondCorp --- Zero Trust de Google}

\begin{concepto}
\textbf{BeyondCorp} (Google, 2014) fue la implementación pionera de Zero
Trust a escala. Después de \textbf{Operación Aurora} (ataque chino de 2010
que comprometió la red interna), Google redesarrolló su modelo de acceso.
\end{concepto}

Principios BeyondCorp:
\begin{enumerate}
    \item \textbf{Device Trust:} cada dispositivo tiene un certificado.
          Los dispositivos no administrados no tienen acceso, sin importar
          dónde estén conectados.
    \item \textbf{User Identity:} autenticación fuerte (MFA) en \textit{cada}
          acceso. La VPN corporativa no existe --- el usuario se autentica
          directamente al recurso.
    \item \textbf{Context-Aware Access:} acceso depende del contexto:
          dispositivo administrado, IP inusual, horario normal, rol correcto.
    \item \textbf{Micro-segmentación:} acceso al sistema de RH no da
          acceso al sistema de pagos.
\end{enumerate}

\begin{lstlisting}
  BeyondCorp Architecture:

  Empleado con laptop corporativa
           │
           ▼
    Access Proxy (BeyondCorp)
           │
           ├── Dispositivo registrado? (certificado)
           ├── Usuario autenticado? (Identity-Aware Proxy)
           ├── Contexto normal? (hora, ubicación, comportamiento)
           └── Tiene permiso para ESTE recurso específico?
                        │
               [TODOS SÍ] → Acceso concedido
               [ALGÚN NO] → Acceso denegado + alerta
\end{lstlisting}

\begin{moderno}
\textbf{\emoji{rocket} Google Cloud BeyondCorp Enterprise} permite que
cualquier empresa implemente Zero Trust usando la misma infraestructura
que protege Google:
\begin{itemize}
    \item \textbf{Identity-Aware Proxy (IAP):} acceso a apps internas sin VPN
    \item \textbf{Endpoint Verification:} verifica el estado del dispositivo
    \item \textbf{Context-Aware Access:} políticas basadas en contexto
    \item \textbf{Chrome Enterprise:} el browser como agente de seguridad
\end{itemize}
Hoy, los 100,000+ empleados de Google trabajan sin VPN corporativa tradicional.
\end{moderno}

\begin{cuidado}
\textbf{Zero Trust no es una tecnología, es una filosofía.} No puedes
``comprar Zero Trust''. Requiere cambios arquitectónicos fundamentales:
inventario de todos los recursos, clasificación de datos, autenticación
fuerte universal, y monitoreo continuo. La implementación típica toma
2--3 años en una empresa grande.
\end{cuidado}
